\subsubsection{Асимптотические обозначения}
\begin{tabular}{ | c | m{15em}| m{15em} | } 
  \hline
  Обозначение & Интуитивное объяснение & Определение \\ 
  \hline
  $f(n) \in O(g(n))$ & $f$ ограничена сверху функцией $g$ (с точностью до постоянного множителя) асимптотически & $\exists C 
 \;\exists N \; : \; \forall n > N \newline f(n) \leqslant C g(n)$ \\ 
  \hline
  $f(n) \in \Omega(g(n))$  & $f$ ограничена снизу функцией $g$ (с точностью до постоянного множителя) асимптотически & $\exists c>0 \;\exists N \; : \; \forall n > N \newline f(n) \geqslant cg(n)$ \\ 
  \hline
  $f(n) \in \Theta(g(n))$ & $f$ ограничена снизу и сверху функцией $g$ асимптотически & $\exists c>0 \; \exists C \;\exists N \; : \; \forall n > N \newline c g(n) \leqslant f(n) \leqslant Cg(n)$ \\ 
  \hline
  $f(n) \in o(g(n))$ & $g$ доминирует над $f$ асимптотически & $\forall c>0 
 \;\exists N \; : \; \forall n > N \newline f(n) \leqslant c g(n)$ \\ 
  \hline
  $f(n) \in w(g(n))$ & $f$ доминирует над $g$ асимптотически & $\forall C \;\exists N \; : \; \forall n > N \newline f(n) \geqslant cg(n)$ \\ 
  \hline
\end{tabular}

\subsubsection{Многоленточная машина Тьюринга, вычисления на ней, измерение времени и памяти. Недетерминированная машина Тьюринга.}
\begin{definition} \textit{Детерминированная $k$-ленточная машина Тьюринга} -- это семёрка
\newline $M = (\Sigma, \Gamma, Q, q_0, \delta, q_{acc}, q_{rej} )$, компоненты которой имеют следующий вид ($\Sigma, \Gamma, Q$ -- конечные непустые):
\begin{itemize}
    \item Конечное множество $\Sigma$ -- входной алфавит.
    \item Другое конечное множество $\Gamma$ -- рабочий алфавит, содержащий все символы, допустимые на лентах, причём $\Sigma \subset \Gamma$, Рабочий алфавит содержит особый символ пробел: $\llcorner\lrcorner \in \Gamma, \llcorner\lrcorner \notin \Sigma$.
    \item Конечное множество $Q$ -- множество состояний, такое что $\Gamma \cap Q = \emptyset$
    \item Есть начальное состояние $q_0 \in Q$.
    \item Функция переходов $\delta \: \ Q\times\Gamma^k \to Q\times(\Gamma\times\{-1, 0, +1\})^k$ определяет поведение машины на каждом шаге. 
    Если машина находится в состоянии $q \in Q$ и обозревает символ $a \in\Gamma$, то $\delta(q, a_1, \ldots , a_k)$ -- это набор вида $(q',a_1', d_1, \ldots, a_k', d_k)$, 
    где $q' \in Q$ -- новое состояние, $a_i'$ -- символ, записываемый на $i$-й ленте вместо $a_i$, и $d_i \in \{-1, 0, +1\}$ -- направление перемещения головки на $i$-й ленте.
    \item Если машина переходит в принимающее состояние $q_{acc} \in Q$ или в отвергающее состояние $q_{rej} \in Q$, то она останавливается.

\end{itemize}
\end{definition}

\begin{definition} 
\textit{Недетерминированная k-ленточная машина Тьюринга} определяется так же, с функцией переходов $\delta \: \ Q\times\Gamma^k \to 2^{Q\times(\Gamma\times\{-1, 0, +1\})^k}$
\end{definition}

\textit{Конфигурация:} состояние, содержимое всех лент, положение всех головок. 

\textit{Вычисление}:
последовательность конфигураций. Для детерминированной машины оно единственно. Для
недетерминированной их может быть более одного.

Вариант: машина с отдельной входной лентой. Первая лента, входная, доступна только
для чтения, на ней записана входная строка $w\in\Sigma^*$, ограниченная меткой начала и
меткой конца. Остальные ленты -- рабочие.

Многоленточная МТ работает за \textbf{время} $t(n)$, если на всяком входе длины $n$ она останавливается не более чем через $t(n)$ шагов. 

Многоленточная МТ с отдельной входной лентой использует \textbf{память} $s(n)$, если на всяком входе длины $n$ ни одна головка на рабочих лентах за время вычисления не отъезжает далее чем на $s(n) - 1$ позицию от своего исходного положения.

\begin{definition} Классом $\dtime{T(n)}$ называется класс языков, которые распознаются за время $O(T(n))$ детерминированной многоленточной машиной Тьюринга (ДМТ). Иными словами, время работы машины на любом слове длины $n$ не превосходит некоторой константы, умноженной на $T(n)$.
\end{definition}

\begin{definition} Классом $\ntime{T(n)}$ называется класс языков, которые распознаются за время $O(T(n))$ недетерминированной машиной Тьюринга (НДМТ).
\end{definition}

Как правило, класс $\dtime{T(n)}$
рассматривают не для всех функций $T$,
а только для ''хороших''. Во-первых, естественно потребовать $T(n) \geqslant n$, иначе машина не успеет даже прочесть свой вход. Во-вторых, естественно потребовать монотонности $T(n)$: чем длиннее вход, тем больше время работы. В-третьих, функция должна быть конструируемой по времени: вычислить $T(n)$
должно быть возможно за время $T(n)$.

\begin{definition}
Функция $T :\ \N \to \N$
называется \textit{конструируемой по времени}, если существует алгоритм, который за время $O(T(n))$
по $1^n$ (т. е. числу $n$ в унарной записи) получает $T(n)$ (в бинарной записи).
\end{definition}

Унарная запись аргумента введена лишь для единообразия: тогда время
вычисления $T(n)$ измеряется как функция от длины аргумента, т.е. $n$,  и это
время не должно быть больше $O(T(n))$. Все ''обычные'' функции конструируемы по времени: $n^c, 2^n, n\log n$ и тд.  Более того, для построения функции, не конструируемой по времени, нужно применить специальную нетривиальную конструкцию. С другой стороны, рассмотрение не конструируемых по времени функций приводит к парадоксальным результатам, например, теорема об
иерархии будет неверна.

\begin{definition}
Функция $T :\ \N \to \N$
называется \textit{конструируемой по памяти}, если существует алгоритм, который
по $1^n$ (т. е. числу $n$ в унарной записи) получает $T(n)$ (в бинарной записи), используя память $O(T(n))$.
\end{definition}

\begin{definition} Классы $\dsp{T(n)}$, $\nsp{T(n)}$ определяются аналогично, но использующие $O(T(n))$ памяти.
\end{definition}

\subsubsection{Временные сложностные классы: \p, \e, \EXP. Замкнутость (или незамкнутость) \p относительно теоретико-множественных
операций.}
\begin{definition} $\p = \bigcup\limits_{c=1}^{\infty}\dtime{n^c} = \dtime{poly(n)}$

Иными словами классом \p называется множество языков, распознающихся за полином от длины входа. Обратите внимание, что если в постановку задачи входит числовой параметр
$n$, то полином считается не от этого параметра, а от его логарифма, т.к.
на запись числа $n$ требуется порядка $\log n$ битов
\end{definition}

\begin{note}
Класс \p замкнут относительно пересечений, объединений и дополнений.
\end{note}

\begin{proof}
Если $A, B \in \p$, то за полиномиальное время мы можем вычислить предикаты $x \in A$ и $x \in B$. Тогда с помощью конъюнкции, дизъюнкции и отрицания можно выразить пересечение, объединение и дополнение.
\end{proof}

\begin{example} Примерами языков из
\p являются такие множества:

\begin{enumerate}
    \item $ \mathsf{GCD} = \{(a,b,d) \ | $ число $d$ является наибольшим общим делителем чисел $a$ и $b$ \}.  Действительно, наибольший общий делитель можно найти алгоритмом Евклида и затем сравнить результат с $d$.
    \item $ \mathsf{PATH} = \{(G,s,t) \ | $ в графе $G$ есть путь из $s$ в $t$ \}. Наличие пути в графе проверяется быстро, например обходом в ширину.
\end{enumerate}
\end{example}

\begin{definition} 
$\e = \bigcup\limits_{c = 1}^\infty \dtime{2^{cn}} = \dtime{2^{O(n)}}$.
\end{definition}

\begin{definition} 
$\EXP = \bigcup\limits_{c = 1}^\infty \dtime{2^{n^c}} = \dtime{2^{poly(n)}}$.
\end{definition}

\subsubsection{Класс \np : определения через сертификаты и через недетерминированные машины.}
\begin{definition} Класс \np состоит в точности из тех языков $A$, для которых существует некая ДМТ двух аргументов $V(x, s)$, работающая по времени за $poly(|x|)$ такая, что
$$
\forall x \in \Sigma^* \ x \in A \Longleftrightarrow \exists s: V(x, s) = 1.
$$
Тогда $V$~--- верификатор, а $s$~--- сертификат для данного $x$.  Таким образом, сертификат удостоверяет, что $x \in A$, а верификатор проверяет верность сертификата. Для слов из языка подходящий сертификат
должен существовать, а для слов не из языка все сертификаты должны отвергаться. Иначе говоря, класс \np -- это класс языков, принадлежность к которым можно быстро доказать, а \p -- класс языков, принадлежность к которым можно быстро выяснить.
\end{definition}

\begin{definition} $\np = \bigcup\limits_{c=1}^{\infty}\ntime{n^c}$
\end{definition}

\begin{note} Определения выше эквивалентны (смотри \ref{th:eqvivNP}).
\end{note} 

\subsubsection{Замкнутость (или незамкнутость) \np относительно теоретико-множественных
операций.}

\begin{statement}
Класс \np замкнут относительно пересечений, объединений и звезды Клини.
\end{statement}

\begin{proof}
Пусть $A, B \in \np$. Тогда по определению \np $x \in A \lra \exists s_1\ V(x, s_1) = 1$ и $x \in B \lra \exists s_2\ W(x, s_2) = 1$. 
Следовательно,
$x \in A \cap B \lra x \in A \wedge x \in B \lra \exists (s_1, s_2) \ V(x, s_1) \wedge W(x, s_2) = 1$.
Аналогично для объединения:
$x \in A \cup B \lra x \in A \lor x \in B \lra \exists (s_1, s_2) \ V(x, s_1) \lor W(x, s_2) = 1$.

Если язык $A$ принадлежит \np, то всякое слово из $A^*$ представляется в виде произведения слов из $A$. Пустые подслова можно не рассматривать (или считать, что такой сомножитель всего один). Длины подслов не превосходят длины слова, равной $O(n)$, и количество подслов также не превосходит этой величины. Угадываем разложение (если оно есть), и для каждого из слов за время $O(n^c)$ получаем сертификат. Тогда за время $O(n^{c+1})$ получается сертификат принадлежности языку $A^*$.
\end{proof}

\begin{note}
Про замкнутость класса \np относительно дополнения ничего не известно.
Если \p = \np, то \np = \conp.
Об истинности обратной теоремы ничего не известно: из \np = \conp автоматически не следует \p = \np.
\end{note}

\begin{proof}
Действительно, пусть $A \in \conp$. Значит, $\overline{A} \in \np$. Поскольку мы предположили, что \p = \np, то $\overline{A} \in \p$. Поскольку класс \p замкнут относительно дополнения, то и $A \in \p$. Значит, $\conp \subset \p$. Поскольку $\p \subset \conp$, имеем $\p = \conp$, а значит и \np = \conp в предположении \p = \np.
\end{proof}

\subsubsection{Принадлежность к \np задач о выполнимости, 3-раскраске, клике, гамильтоновом пути, рюкзаке и других подобных.}

\begin{example} Примерами языков из
\np являются такие множества:

\begin{enumerate}
    \item $ \mathsf{SAT} = \{\phi \ | $ $\phi$ выполнимая булева формула \}.   (Выполнимость означает, что формула равна 1 на некотором наборе значений). В данном случае сертификатом будет выполняющий набор, а верификатор проверит, что значение формулы на этом наборе действительно равно 1. Важным частным случаем является задача $\mathsf{3SAT}$, в которой про формулу $\phi$ дополнительно сказано, что она представлена в виде 3-КНФ, т.е. КНФ, в которой каждый дизъюнкт включает в себя 3 литерала.
    \item $ \mathsf{CLIQUE} = \{(G,k) \ | $ в графе $G$ найдётся полный подграф хотя бы на $k$ вершинах \}. Сертификатом будет сама клика.
    \item $ \mathsf{3COL} = \{G \ | $ вершины графа $G$ можно правильно раскрасить в 3 цвета \}. Сертификатом будет раскраска.
    \item $ \mathsf{HAMCYCLE} = \{G \ | $ в графе $G$ существует гамильтонов цикл \}. Сертификатом будет сам цикл.
    \item $ \mathsf{KNAPSACK} = \{ \ (n_1,\ldots, n_k; m_1, \ldots m_k; N; M) \ | \ \exists \alpha \in \{0,1\}^k \sum{\alpha_i n_i} \leq N$ и $\sum{\alpha_i m_i} \geq M$  \}. Сертификатом будет набор $\alpha$.
\end{enumerate}
\end{example}

\subsubsection{Сводимость задачи о выполнимости произвольной формулы к задаче о выполнимости 3-КНФ.}

\begin{statement}
Докажем, что $\mathsf{SAT} \karp \mathsf{3SAT}$
\end{statement}

\begin{proof}
Пусть $\phi$ -- пропозициональная формула. Введём переменную для каждой её подформулы (включая исходные переменные и всю формулу).
Каждая подформула, составленная из двух, задаёт утверждение вида $q = r \wedge s$. Его можно представить как 3-КНФ (в данном случае $(q \vee \neg r \vee \neg s) \wedge (\neg q \vee r \vee \neg s) \wedge (\neg q \vee \neg r \vee s) \wedge (\neg q \vee r \vee s)$). Взяв конъюнкцию всех таких формул, мы получим 3-КНФ, выполнимость которой эквивалентна выполнимости исходной формулы. Действительно, выполняющий набор исходной формулы задаст значения всех подформул, которые будут выполняющим набором полученной формулы. И наоборот, из выполняющего набора новой формулы можно выделить выполняющий набор исходной.

Осталось пояснить, почему этот алгоритм работает полиномиальное время.
Ясно, что достаточно построить за полиномиальное время дерево синтаксического разбора формулы, дальнейшие операции занимают константное время
для каждой его вершины. Построение дерева также несложно: путём подсчёта
скобочного итога можно выделить из формулы две подформулы, из которых
она получена, и повторить процедуру с ними рекурсивно. Это займёт линейное
время для каждой подформулы, т.е. всего квадратичное время.
\end{proof}

\subsubsection{Сводимость задачи о 3-раскраске к задаче о выполнимости.}
\begin{statement}
Любую задачу можно свести к задаче о выполнимости. Для примера покажем, как это делается для $\mathsf{3COL}$.
\newline То есть докажем, что $\mathsf{3COL} \karp \mathsf{SAT}$
\end{statement}

\begin{proof}
Заведём для каждой вершины $v$ две переменные: $p_{v1}$ и $p_{v2}$.
Наложим на них условие $p_{v1} \vee p_{v2}$. Это значит, что для каждой вершины есть
три варианта логических значений этой пары: $01, 10$ и $11$. Осталось добавить
условие, что для соседних вершин эти варианты разные: если есть ребро $(v, w)$,
то добавляем условие $(p_{v1} \neq p_{w1}) \vee (p_{v2} \neq p_{w2})$. Итоговая формула будет конъюнкцией всех условий. Очевидно, любая правильная раскраска перекодируется
в выполняющий набор, и наоборот.
\end{proof}

\subsubsection{Сводимость задач о клике, независимом множестве и вершинном покрытии друг к другу.}
Смотри билет \ref{sec:question4}

\subsubsection{Сводимость задачи о гамильтоновом пути к задаче о гамильтоновом цикле.}
Смотри билет \ref{sec:question7}

\subsubsection{Класс \conp : определения через сертификаты и через недетерминированные машины.}

\begin{definition}
Класс \conp состоит в точности из тех языков $A$, для которых $\overline{A} \in \np$. \newline То есть $\conp = \{A \; : \: \overline{A} \in \np\}$
\newline Или альтернативно: $\forall x \in \Sigma^* \ x \in A \Longleftrightarrow \forall s: V(x, s) = 1.
$ 
\end{definition}

Иначе говоря, если к \np
относятся языки, принадлежность к которым легко доказать, то \conp -- класс языков, принадлежность к которым легко опровергнуть.

\begin{example}
Самым естественным представителем класса
\conp является язык
тавтологий $\mathsf{TAUT} = \{\phi\ | \ \phi$ -- тавтология \}.  Действительно, опровергнуть, что $\phi$
тавтология, очень легко: нужно предъявить набор значений, на котором формула ложна. Доказать, что $\phi$
тавтология, можно, предъявив вывод в исчислении высказываний, однако такой вывод может быть слишком длинным: длиннее,
чем полином от длины $\phi$. Коротких доказательств тавтологичности не известно.
\end{example}

\subsubsection{Полиномиальная сводимость по Карпу. Свойства.}
\begin{definition}
Пусть $A$ и $B$ суть два языка. Тогда
$A$ \textit{сводится по Карпу}
к $B$, если существует всюду определённая функция $f \; : \{0, 1\}^* \to \{0, 1\}^*$, вычислимая за полиномиальное время, такая что
$x\in A \; \Leftrightarrow \; f(x)\in B$. Обозначение: $A \karp B$. (Индекс $p$
означает полиномиальность).
\end{definition}

\begin{statement}
    Имеют место следующие факты:

    \begin{enumerate}
        \item Полиномиальная сводимость рефлексивна: $A \karp A$;
        \item Полиномиальная сводимость транзитивна: если $A\karp B$ и $B \karp C$, то $A \karp C$;
        \item Если $A\in\p$, а $B\neq \emptyset$ и $B\neq\{0,1\}^*$, то $A \karp B$;
        \item Если $B\in \p$ и $A\karp B$, то $A\in \p$;
        \item Если $B\in \np$ и $A\karp B$, то $A\in \np$;
    \end{enumerate}
\end{statement}

\begin{proof} Пойдем по порядку:
\begin{enumerate}
    \item Рефлексивность очевидна: достаточно рассмотреть функцию $f(x)=x$
    \item  Для транзитивности нужно рассмотреть композицию: если $f$ сводит
$A$ к $B$, а $g$ сводит $B$ к $C$, то $g \circ f$ сводит $A$ к $C$: $x \in A \;\Leftrightarrow \; f(x)\in B \;\Leftrightarrow \; g(f(x)) \in C$.
При этом если и $f$, и $g$ вычисляются за полиномиальное время, то $g \circ f$ также
вычисляется за полиномиальное время.
\item Можно
зафиксировать $b_1\in B$ и $b_2 \notin B$ и рассмотреть функцию  $ f(x) =
  \begin{cases}
    b_1,    & \;\; x \in A\\
    b_2.  & \;\; x \notin A
  \end{cases}
$ \newline В силу разрешимости $A$ эта функция будет вычислимой. Впрочем: тут имеет место неконструктивность: откуда брать нужные $b_1$ и $b_2$, неясно. Ведь множество
$B$ может быть даже неразрешимым.
\item Четвёртое утверждение самое важное и полезное, но также несложно. Достаточно заметить, что индикаторная функция $I_A$ представляется как
композиция $I_B \circ f$. Если $B \in \p$, то $I_B$ вычисляется за полиномиальное время,
а если $A \karp B$, то соответствующая $f$ также вычисляется за полиномиальное
время. Значит, их композиция, т.е. $I_A$, тоже вычисляется за полиномиальное
время, что и означает $A\in\p$.
\item Докажем двумя способами. 
\begin{itemize}
    \item Во-первых, для
определения через недетерминированные машины достаточно заметить, что если $I_B$ вычисляется недетерминированной машиной за полиномиальное время,
а $f$(детерминированно) вычисляется за полиномиальное время, то $I_A = I_B \circ f$
также вычисляется недетерминированной машиной за полиномиальное время.
\item Во-вторых, для определения через сертификаты: если $W(y,s)$ будет верификатором принадлежности $y$ к $B$, то $V(x,s) = W(f(x),s)$ будет верификатором
принадлежности $x$ к $A$: если $x\in A$, то для $f(x)$ будет существовать сертификат и $V$ его примет, а если $x\notin A$, то для $f(x)$ сертификата существовать не
будет.
\end{itemize}


\end{enumerate}

\end{proof}

\subsubsection{\nph и \npc.}

\begin{definition}
Язык $B$ является \nph, если для любого $A \in \np$
выполнено $A \leqslant_p B$.
\end{definition}
\begin{definition}
Язык $B$ является \npc, если он \nph и лежит в \np.
\end{definition}

\begin{statement}
Если $A$ является \nph и $A\leqslant_p B$, то $B$ тоже является \nph
\end{statement}

\begin{proof}
Любой язык $C$ из \np
сводится к $A$, а по
транзитивности и к $B$.
\end{proof}

\subsubsection{Теорема об иерархии по времени}

Один из инструментов, позволяющий доказать,
например, $\p \subsetneq \e \subsetneq \EXP$ -- это теорема об иерархии по времени. Неформально говоря, теорема об иерархии гласит, что если функция $g(n)$ растёт существенно быстрее, чем $f(n)$, то за время $g(n)$ можно распознать больше языков, чем за
время $f(n)$. Точная формулировка такая:

\begin{theorem}
(об иерархии по времени). Пусть $f : \N \to \N$ и $g : \N \to \N$ -- строго
монотонные конструируемые по времени функции, такие что $f(n) \log f(n) = o(g(n))$. Тогда $\dtime{f(n)} \subsetneq \dtime{g(n)}$.
\end{theorem}

\begin{proof}
    Смотри \ref{th:hir-time}
\end{proof}

\begin{corollary}
Существуют строго монотонные функции, не конструируемые
по времени
\end{corollary}

\begin{proof}
    Смотреть в \ref{notconstr}
\end{proof}

\subsubsection{Уровни полиномиальной иерархии.}

\begin{definition}
Классом \textbf{$\Sigma_k^p$}
называется множество языков $A$, для которых
существует полиномиально вычислимый предикат $V$, такой что
    $$x \in A \; \Leftrightarrow \; \exists y_1 \forall y_2 \exists y_3 \ldots (\exists / \forall) y_n \ V(x, y_1,\ldots ,x_n) = 1$$
(последний квантор будет $\forall$ при чётном $n$ и $\exists$ при нечётном).   Аналогично классом \textbf{$\Pi_k^p$}
называется множество языков $A$, для которых
существует полиномиально вычислимый предикат $V$, такой что
    $$x \in A \; \Leftrightarrow \; \forall y_1 \exists y_2 \forall y_3 \ldots (\exists / \forall) y_n \ V(x, y_1,\ldots ,x_n) = 1$$
(последний квантор будет $\exists$ при чётном $n$ и $\forall$ при нечётном).
Классом \ph
называется объединение всех \textbf{$\Sigma_k^p$}.
\end{definition}

\begin{note}
Из определения сразу видно, что $\Sigma_0^p = \Pi_0^p = \p$, $\Sigma_1^p = \np$, $\Pi_1^p = \conp$. 
\end{note}

\subsubsection{Вложенность классов полиномиальной иерархии.}

\begin{statement}
Язык $A$ лежит в $\Sigma_k^p$ тогда и только тогда, когда $\overline{A}$ лежит в $\Pi_k^p$.
\end{statement}

\begin{proof}
\begin{align*}
    x & \in \overline{A} \Leftrightarrow \; x \notin A \\
    \Leftrightarrow \; & \neg \exists y_1 \forall y_2 \exists y_3 \ldots (\exists / \forall) y_n \ V(x, y_1,\ldots ,x_n) = 1 \\
    \Leftrightarrow \; & \forall y_1 \exists y_2 \forall y_3 \ldots (\exists / \forall) y_n \ V(x, y_1,\ldots ,x_n) = 0 
\end{align*}
 Заменив предикат $V$ на его отрицание, получаем формулу для $\Pi_k^p$, что и требовалось.
\end{proof}

\begin{statement}
    $\Sigma_k^p \cup \Pi_{k}^p \subset \Sigma_{k+1}^p \cap \Pi_{k+1}^p$
\end{statement}

\begin{proof}
Нужно доказать 4 отдельных вложения каждого из классов
слева в каждый из классов справа. Каждое из них доказывается добавлением
фиктивной переменной в предикат $V$ и квантора по ней в начало или в конец
цепочки. Например, для доказательства вложения $\Sigma_k^p \subset \Pi_{k+1}^p$ нужно добавить
квантор всеобщности по новой переменной $y_0$ в начало.
\end{proof}

\subsubsection{Замкнутость (или незамкнутость) классов полиномиальной иерархии относительно теоретико-множественных операций.}

\begin{statement}
Если $A \in \Sigma_k^p$ и $B \in \Sigma_k^p$, то $A\cup B \in  \Sigma_k^p$, $A\cap B \in \Sigma_k^p$. Аналогично для $\Pi_k^p$.

Если $A \in \Sigma_k^p$ и $B \in \Pi_k^p$, $A\cup B$ и $A\cap B$ лежат в $\Sigma_{k+1}^p \cap \Pi_{k+1}^p$.
\end{statement}

\begin{proof}
Все эти утверждения следуют из процедуры приведения формулы к предварённой нормальной форме. Например, пусть
$$x \in A \; \Leftrightarrow \; \exists y_1 \forall y_2 \exists y_3 \ldots (\exists / \forall) y_n \ V(x, y_1,\ldots ,x_n) = 1$$
и
$$x \in B \; \Leftrightarrow \; \exists y_1 \forall y_2 \exists y_3 \ldots (\exists / \forall) y_n \ W(x, y_1,\ldots ,x_n) = 1$$
Тогда
\begin{align*}
    x & \in A\cap B \\
    \Leftrightarrow \; & \exists y_1 \forall y_2 \exists y_3 \ldots (\exists / \forall) y_n \ V(x, y_1,\ldots ,x_n) = 1 \wedge \exists z_1 \forall z_2 \exists z_3 \ldots (\exists / \forall) z_n \ W(x, z_1,\ldots ,z_n) = 1 \\
    \Leftrightarrow \; & \exists y_1 \exists z_1 \forall y_2 \forall z_2 \exists y_3 \exists z_3 \ldots (\exists / \forall) y_n (\exists / \forall) z_n \ \left(V\left(x, y_1,\ldots ,x_n\right) = 1 \wedge  W\left(x, z_1,\ldots ,z_n\right) = 1 \right) \\
    \Leftrightarrow \; & \exists (y_1, z_1) \forall (y_2, z_2) \exists (y_3, z_3) \ldots (\exists / \forall) (y_n, z_n) \ \left(V\left(x, y_1,\ldots ,x_n\right) = 1 \wedge  W\left(x, z_1,\ldots ,z_n\right) = 1 \right)
\end{align*}
Последняя запись показывает принадлежность к $\Sigma_k^p$, поскольку стоящий после
всех кванторов предикат вычислим за полиномиальное время, если использовать полиномиально вычислимое кодирование пар.

Аналогично доказываются и остальные включения. Принадлежность к $\Sigma_{k+1}^p \cap \Pi_{k+1}^p$ за счёт двух возможных порядков приведения к предварённой нормальной
форме.
\end{proof}

\begin{corollary}
    $\mathbf{PH} = \cup_{k=1}^{\infty}\Pi_{k}^p$
\end{corollary}

\subsubsection{Классификация по уровням полиномиальной иерархии задач о минимальной эквивалентной формуле, об обобщённой рамсеевости, о кликовой раскраске. Пример задачи на третьем или более высоком уровне полиномиальной иерархии.}
Помимо задач оптимизации есть и другие естественные примеры задач с
различных уровней иерархии (обычно второго или третьего), например:

\begin{itemize}
    %\item \textit{Задача о минимальной ДНФ.} Рассматривается язык $\mathsf{MINEQDNF}  = \{ (\phi, k) \ | $ у формулы $\phi$ есть эквивалентная ей ДНФ длины не более $k \}$. Он принадлежит $\Sigma_k^p$, поскольку задаётся формулой $\exists \psi \forall x \; (\phi(x) = \psi(x) \vee |\psi| \leq k )$. Заметим, что соответствующая задача точной оптимизации будет лежать в $\Sigma_{3}^p \cap \Pi_{3}^p$.

    \item \textit{Задача о минимальной эквивалентной формуле.} Рассмотрим язык  $\{ (\phi, k) \ | $ у формулы $\phi$ есть эквивалентная ей формула $\psi$ длины не более $k \}$. Он принадлежит $\Sigma_2^p$, поскольку задаётся формулой $\exists \psi \forall x \; (\phi(x) = \psi(x) \vee |\psi| \leq k )$. Проверка условия полиномиальна.
    
    \item \textit{Задача о кликовой раскраске}. Пусть дан граф $G = (V, E)$. Его кликовой раскраской в $k$ цветов называется такая функция $f : V \to \{ 1, \ldots, k \}$, что любая максимальная клика в графе содержит вершины хотя бы двух разных цветов. Язык $\{ (G, k) \ | \ $ у графа $G$ есть кликовая раскраска в $k$ цветов $\}$ лежит в $\Sigma_{2}^p$: формула имеет вид $\exists f \forall X \subset V$ (если $X$ — максимальная клика, то $f(X)$ cодержит хотя бы два элемента). Действительно, удостовериться, что $X$ клика, можно путём проверки, что все пары вершин из $X$ образуют рёбра. Проверить максимальность тоже легко: любая другая вершина должна быть не соединена хотя бы с одной из вершин $X$ Отметим, что в отличие от обычной $k$-раскрашиваемости, кликовая $k$-раскрашиваемость может не сохраниться при переходе к подграфу. Это происходит от того, что набор максимальных клик может измениться. Например, полный граф на 5 вершинах можно кликово раскрасить в 2 цвета, а цикл из 5 вершин уже нельзя. Если потребовать, чтобы все подграфы также допускали кликовую раскраску в $k$ цветов, то соответствующий язык будет лежать в $\Pi_{3}^p$. 

    \item \textit{Задача о рамсеевости частичной раскраски рёбер графа.} Известна теорема Рамсея: для любых натуральных $k$ и $l$ при всех достаточно больших $n$ в любой раскраске рёбер полного графа на $n$ вершинах в красный и синий цвета найдётся либо красная клика на $k$ вершинах, либо синяя клика на $l$ вершинах. Более общая задача ставится так: даны три графа, $F$, $G$ и $H$. Верно ли, что для любой раскраски рёбер $F$ в красный и синий цвета найдётся либо красный подграф, изоморфный $G$, либо синий подграф, изоморфный $H$? Соответствующий язык будет лежать в $\Pi_2^p$, формула будет иметь вид $\forall f : E(F) \to \{r, b\} \ \exists K ((K \cong G \wedge f |_{E(K)} \equiv r) \vee (K \cong H \wedge f |_{E(K)} \equiv b))$, здесь $E(F)$ обозначает множество рёбер графа $F$, $f$ -- раскраска (отображение в множество из двух цветов),$f |_{E(K)}$ -- ограничение раскраски на подграф $K$.
    
    % \item \textit{Задача о размерности Вапника–Червоненкиса.} Размерностью Вапника–Червоненкиса $VC(\mathcal{S})$ системы множеств $\mathcal{S} = (S_1, \ldots, S_k)$ называется максимальное количество элементов в множестве $X$, таком что среди множеств $S_i \cap X$ встречаются все подмножества $X$. Предположим, что $\mathcal{S}$ задана в сжатой форме, например, формулой $\phi$ полиномиальной длины, такой что для $i \in \{1, \ldots, k \}$ и $y$ истинность $\phi(i, y)$ эквивалентна принадлежности $y$ к $S_i$. Соответствующую систему обозначим через $\mathcal{S}_\phi$. (Можно заменить формулу на схему из функциональных элементов или другое короткое эффективно вычислимое описание). Тогда язык $\mathsf{VCDIM} = \{(\phi, m) \ | \ VC(\mathcal{S}_\phi) \geq m\}$ лежит в $\Sigma_{3}^p$. Действительно, должно найтись такое $X$ из $k$ элементов, что для любого $A \subset X$ найдётся $i$, такое что $A = \{ y\in X \ | \ \phi(i, y) = 1 \}$. Если же стоит задача точного подсчёта размерности, то соответствующий язык будет лежать в $\Sigma_{4}^p \cap \Pi_{4}^p$.
\end{itemize}


\subsubsection{Классы $\mathbf{PSPACE}$, $\mathbf{NPSPACE}$. Доказательство $\mathbf{PH} \subset \psp \subset \EXP$. Вывод $\mathbf{PSPACE} = \mathbf{NPSPACE}$ из теоремы Сэвича.}

\lsp, \nlsp.
\begin{definition} 
$\mathbf{PSPACE}=\bigcup_{c=0}^{\infty}\dsp{n^c}$
\end{definition}

\begin{lemma}
    Класс $\mathbf{PSPACE}$ замкнут  относительно теоретико-множественных операций.
\end{lemma}

\begin{proof}
    Если операция отрицание, то делаем инвертирование входа. Если бинарная операция, то отдельно на каждую программу выделяем по $poly$ памяти, а потом просто булева функция, чтобы ответы склеить. $poly + poly = 2 poly$, что всё ещё $poly$.
\end{proof}

\begin{definition} 
$\mathbf{NPSPACE}=\bigcup_{c=0}^{\infty}\nsp{n^c}$
\end{definition}

\begin{theorem}
Имеют место соотношения 
\begin{align*}
    \mathbf{DTIME}(t(n)) \subset \mathbf{NTIME}(t(n)) \subset \dsp{t(n)} \subset \nsp{t(n)} \subset \\ \subset \mathbf{DTIME}(2^{O(t(n))}) = \bigcup_{c=0}^{\infty}\mathbf{DTIME}(2^{ct(n)})
\end{align*}
\end{theorem}

\begin{corollary}
    $\p \subset \np \subset \mathbf{PSPACE} \subset \mathbf{NPSPACE} \subset \EXP$ и $\mathbf{PH} \subset \psp \subset \EXP$
\end{corollary}

\begin{proof}
    Первое теормин не просит, поэтому докажем второе:
    
Про $\mathbf{PH} \subset \psp$: Возьмем язык $A \in \mathbf{PH}$. Он задается машиной Тьюринга $M$ с формулой с переменным количеством кванторов. Пусть наша полиномиальная по памяти машина делает следующее: выделит на каждую кванторную переменную полиномиальное количество места и столько на работу $M$. И пусть тупо перебирает все биты, а потом запускает $M$ на всех этих входах.

Про $\EXP$: т.к. всего ячеек полиномиально, то всего конфигураций $\Gamma^{poly(n)}$, т.е. экспоненциально. Это значит, что за экспоненциальное время у нас либо машина остановится, либо у неё повторится конфигурация, т.е. она зациклится.
\end{proof}

\begin{theorem}
    (Теорема об иерархии по памяти)
    Пусть $f$ и $g$ -- возрастающие функции, конструируемые по памяти, причём $f(n) = o(g(n))$ и $f(n) \geq \log n$. Тогда $\dsp{f(n)}\subsetneq \dsp{g(n)}$ Аналогичные вложения верны и для $\mathbf{SPACE}$ и $\mathbf{NSPACE}$.
\end{theorem}

\begin{proof}
    См. задачу \ref{task:memory}.
\end{proof}

\begin{theorem}[Сэвич] Если $s(n) \ge \log n$, то $\nsp{s(n)} \subset \dsp{s(n)^2}$. Более того, $\psp = \npsp$.
\end{theorem}

\begin{proof}
    См. билет \ref{sec:savitch}
\end{proof}

\subsubsection{Принадлежность к $\mathbf{PSPACE}$ задачи о булевых формулах с кванторами и задачи
о выигрышной стратегии в обобщённой игре в города.
}

Как уже говорилось, мы не умеем доказывать $\np \neq \psp$ и даже
$\p \neq \psp$. Однако, как и в классе $\np$, в классе $\psp$ есть ''самые сложные'' задачи, которые называются $\mathbf{PSPACE}$-полными. Общая идея такая же как и раньше: язык должен лежать в $\psp$, а любой другой язык из
$\psp$ должен к нему сводиться. Вопрос лишь в том, какую сводимость
использовать: не нужно ли наложить на сводящую функцию условие полиномиальности по памяти? Разумеется, не нужно: нас же интересует отделимость
$\psp$ от $\p$, поэтому сводимость должна быть из более маленького класса.
Поэтому оставим сводимость по Карпу.

\begin{definition}
    Язык $B$ называется $\mathbf{PSPACE}$-трудным, если для любого языка $A \in \mathbf{PSPACE}$ выполнено $A \karp B$. 
\end{definition}

\begin{definition}
    Язык называется $\mathbf{PSPACE}$-полным, если он $\mathbf{PSPACE}$-труден и лежит в $\mathbf{PSPACE}$.
\end{definition}

Как и прежде, выполнено простое утверждение

\begin{statement}
Если $A$ является $\mathbf{PSPACE}$-трудным и $A\leqslant_p B$, то $B$ тоже является $\mathbf{PSPACE}$-трудным
\end{statement}

 $\mathbf{PSPACE}$-полных задач известно не так много, как \npc, но всё же существенное количество. 
 
 Рассматривая теорию \npc, мы доказали полноту задачи $\sat$, а затем уже её сводили к другим. В теории $\mathbf{PSPACE}$-полноты тоже будет такая задача: $\mathsf{TQBF}$.

\begin{definition}
    Языком $\mathsf{TQBF}$ называется множество булевых формул $\phi$, таких что для некоторого $x_1 \in \{ 0, 1 \}$ найдётся $x_2 \in \{ 0, 1 \}$, такое что для некоторого $x_3 \in \{ 0, 1 \}$ найдётся \ldots (цепочка чередующихся кванторов по всем переменным) $\phi(x_1, x_2, x_3, \ldots)$ истинна.

\end{definition}

Буквы $\mathsf{TQBF}$ означают ''totally quantified boolean formulae'', по-русски эту задачу называют задачей о булевых формулах с кванторами и иногда обозначают
БФК. Заметим, что требование строгого чередования кванторов излишне: если изначально два одноимённых квантора идут подряд, то можно между ними
вставить квантор по новой переменной, от которой $\phi$ на самом деле не зависит. (Если требовать обязательного вхождения новой переменной, можно взять
конъюнкцию $\phi$ с $z \vee \neg z$. Формально можно сказать, что язык $\mathsf{TQBF}'$
, в котором
кванторы могут не чередоваться, сводится к обычному $\mathsf{TQBF}$.

\begin{theorem}
     Язык $\mathsf{TQBF}$ является $\mathsf{PSPACE}$-полным.
\end{theorem}

\begin{proof}
    См билет \ref{sec:pspc}
\end{proof}

\begin{theorem}
     Принадлежность к $\mathbf{PSPACE}$ задачи о выигрышной стратегии в обобщённой игре в города.
\end{theorem}

\begin{proof}
    См билет \ref{sec:pspc2}
\end{proof}


\subsubsection{Классы: $\mathbf{L}$, $\mathbf{NL}$. Верны вложения: $\mathbf{L} \subset \p$, $\mathbf{NL} \subset \p$ }
 

\begin{definition}
$\mathbf{L} = \dsp{\log n}$.
\end{definition}

\begin{definition}
$\mathbf{NL}=\nsp{\log n}$
\end{definition}

\begin{example}
    Примеры языков
    \begin{itemize}
        \item Язык $\mathsf{LE} = \{(x, y) \ | \ x \leq y\}$ принадлежит $\mathbf{L}$. (Имеется в виду, что числа $x$ и $y$ записаны в двоичной записи, возможно, с ведущими нулями.)

        \item Язык $\mathsf{ADD} = \{(x, y, z) \ | \ x + y = z \}$ принадлежит $\mathbf{L}$.

        \item Язык $\mathsf{MULT} = \{(x, y, z) \ | \ x \cdot y = z \}$ принадлежит $\mathbf{L}$.

        \item Язык правильных скобочных последовательностей для нескольких типов скобок лежит в $\mathbf{L}$.

        \item Язык $\mathsf{TREE} = \{ G \ |$ неориентированный граф $G$ является деревом $\}$ лежит в $\mathbf{L}$.
    \end{itemize}
\end{example}

\begin{note}
$\mathbf{L} \subseteq \mathbf{NL} \subseteq \p.$ 

%NOOO, GOD, PLEASE, NOOOO!!
% Мусатов He\mathbf{L}\mathbf{P}
%Но $\mathbf{L} \neq \p$. Первое вложение тривиально, второе следует из теоремы о связях
%пространственно-временных классов (теорема 6.5 из compl-book)
\end{note}


\begin{proof}
$\mathbf{L} \subseteq \mathbf{NL}$ - тривиально.

$\mathbf{NL} \subseteq \p$.
  Пусть $A \in \mathbf{NL}$. Недетерминированная машина Тьюринга $M$, распознающая $A$, тратит $\leq C\log(n)$ памяти. Тогда конфгураций у этой машины не более $\Gamma^{C\log(n)} \sim Dn^C$, т.е. полиномиальное количество. Представив конфигурации в виде  ориентированного графа, где ребро означает возможность перейти по функции перехода из одной конфигурации в другую. Тогда положительный ответ машины $M$ равносилен наличию пути от стартовой конфигурации до конфигурации с состоянием $q_{\text{accept}}$. Наличие пути в графе полиномиального размера мы умеем проверять за полином, т.е. такая задача лежит в $\p$.  
\end{proof}

\begin{corollary}
    $\mathbf{L}$ и $\mathbf{NL}$ замкнуты относительно теоретико-множественных операций (пересечения, объединения и дополнения).
\end{corollary}

\begin{proof}
    Если бинарная операция, то отдельно на каждую программу выделяем по $\log$ памяти, а потом просто булева функция, чтобы ответы склеить. $\log + \log = 2 \log$, что всё ещё $\log$.

    Если операция отрицание, то делаем инвертирование входа. 
    
    Но вот с $\mathbf{NL}$ (и вообще N-что-угодно) говорить про замкнутость относительно дополнения, так как выше не выйдет. Причина: нельзя вызывать из недетерминистической машины Тьюринга как подпрограмму недетерминистическую машину Тьюринга, так как будут беды. Нужно другое док-во.
\end{proof}

\subsubsection{Принадлежность к $\mathbf{L}$ задач о сравнении и сложении двоичных чисел, о палиндроме, о сбалансированной последовательности скобок одного вида.}
\begin{statement}
    Язык $\mathsf{LE} = \{(x, y) | x \leq y\}$ принадлежит \lsp. (Имеется в виду, что числа $x$ и $y$ записаны в двоичной записи, возможно, с ведущими нулями.)
\end{statement}

\begin{proof}
    Вначале нужно сравнить количество битов в $x$ и $y$. Для этого нужно пройтись по записи $x$ до первой единицы, завести счётчик и увеличивать его на единицу при чтении нового бита, пока запись не кончится. Затем нужно сделать то же самое для $y$. Полученные счётчики нужно сравнить. Поскольку они имеют логарифмическую длину, их можно сравнить совсем простым алгоритмом, например, вычитая по единице из каждого счётчика, пока один не обратится в ноль. Можно запустить и рекурсию. 
    
    Теперь, если в записи одного из чисел битов больше, то оно и будет больше. Иначе нужно поочерёдно сравнивать соответствующие биты, пока один не будет больше. Тогда число, из которого получен этот бит, будет больше. Если же все биты совпали, значит, и числа равны.
\end{proof}

Задачи о сложении, умножении и ПСП см. в билете \ref{sec:arithmetic}
\begin{statement}
    Задача о палиндроме лежит в \lsp
\end{statement}

\begin{proof}
    Идея такая же как в предыдущих задачах: храним рассматриваемую позицию (как в сложении), идем с начала и с конца и сравниваем символы
\end{proof}

\subsubsection{Принадлежность к $\mathbf{NL}$ задач о пути в ориентированном графе и о сильной связности.}

См. билет \ref{sec:sconnected}

\subsubsection{Логарифмическая сводимость.}

Часто в теории сложности мы задаемся вопросом $A = B$, где $A$ и $B$ это два класса, причем $A \subset B$. Например, $\p = \np$ один из таких вопросов. В попытке доказать это утверждение мы обозначили подкласс $\npc$ задач и сказали: "Если хоть одна из этих задач принадлежит $\p$, то $\p = \np$. Очень важным инструментов здесь предстала \emph{полиномиальная сводимость} $\leq_P$. 

Теперь зададимся похожим вопросом, но для $\mathbf{L} \subset \p$, или $\mathbf{L} \subset \mathbf{NL}$. Однако сводимость по Карпу $\leq_P$ здесь уже не поможет - все задачи, хоть из $\mathbf{L}$, из $\mathbf{NL}$, из $\p$ решаются полиномиально. Потребуется более тонкая, \emph{логарифмическая сводимость} $\leq_l$.

\begin{definition}
Функция $f$ называется логарифмически вычислимой, если существует машина Тьюринга, работающая на входе $x$ за $O(\log(|x|))$ дополнительной памяти, причем на входной ленте можно только читать, а на выходной - только писать.  
\end{definition}

\begin{lemmanote}
Ограничение на чтение, письмо существенно, так как сам вход занимает порядка логарифма памяти, а выход может существенно превышать этот объём. 
\end{lemmanote}



\begin{lemma} Следующие условия эквивалентны:

1. $f$ логарифмически вычислима.

2. Языки $\{ (x, n) \mid |f(x)| \leq n \}$ и $\{(x, j) \mid f(x)|_j = 1\}$, при этом максимальное $n$, при котором пара $(x, n)$ входит в первый язык, ограничено $poly(|x|)$, лежат в $\mathbf{L}$.
\end{lemma}

\begin{proof}

    Смотри compl-book. Страница 120.

    %$1 \rightarrow 2$. Тривиально. Вычисляем $f(x)$ и проверяем условия.
    
    %$2 \rightarrow 1$. Программа (индексация с 1):
\begin{lstlisting}[escapeinside={(*}{*)}]
input(x)
n = 1
while True:
   if |f(x)| >= n:
       if f(x)[n] = 1:
           print(1, end='')
       else: 
           print(0, end='')
       n += 1
   else: break
\end{lstlisting}

%$n$ занимает $\log(|x|)$ памяти, на каждое обращение к данным по условию языкам в $if$-ах тоже порядка логарифма памяти. Причем условия на вход/выход выполняются. Итого: $f$ логарифмически вычислима. 
\end{proof}

\subsubsection{\nlsp-полнота.}
\begin{definition}
    Язык $B$ называется NL-трудным, если для любого $A \in \nlsp$ выполнена сводимость $A \leq_l B$. Язык называется \nlsp-полным, если он \nlspтруден и лежит в \nlsp.
\end{definition} 

Как всегда, выполнены простые утверждения:
\begin{statement}
    Если $B$ является \nlsp-трудным и $B \leq_l C$, то $C$ также
\nlsp-труден.
\end{statement}
\begin{statement}
    Если $B$ является \nlsp-трудным и лежит в \lsp, то \lsp =
\nlsp.
\end{statement}
\subsubsection{Схемы из функциональных элементов}



\subsubsection{$\ppoly$, $\mathbf{NC}^d$, $\mathbf{AC}^d$}

\begin{definition}
    Классом $\mathbf{SIZE}(s(n))$ называется множество языков $A$, для которых существует семейство схем $\{C_n\}_{n=1}^{\infty}$ размера $O(s(n))$, такое что $x \in A$ тогда и только тогда, когда $C_{|x|}(x) = 1$.
\end{definition}

\begin{definition}
    $\ppoly = \bigcup_{c=1}^\infty \mathbf{SIZE}(n^c)$ 
\end{definition}

\begin{definition}
    Классом $\mathbf{NC}^d$ называется класс языков, которые распознаются семейством схем полиномиального размера и глубины $O(\log^d n)$, в которых элементы $\wedge$ и $\vee$ имеют входную степень ровно $2$.
\end{definition}

\begin{definition}
     Классом $\mathbf{NC}$ называется объединение $\mathbf{NC}^d$ для всех натуральных $d$.
\end{definition}

\begin{definition}
    Классом $\mathbf{AC}^d$ называется класс языков, которые распознаются семейством схема полиномиального размера и глубины $O(\log^d n)$, в которых элементы $\wedge$ и $\vee$ имеют любую входную степень.
\end{definition}

\subsubsection{$\p \subsetneq \ppoly$}

\begin{proof}
\begin{itemize}

\item $\p \subset \ppoly$

    Доказательство в целом повторяет доказательство теоремы
Кука–Левина. А именно, если есть машина Тьюринга, работающая за $p(n)$ шагов, то строится схема размера $O(p(n)^2)$, моделирующая работу этой машины.
Схема будет вычислять символы, стоящие в протоколе работы этой машины. Поскольку каждый символ зависит только от 4 символов на предыдущем этапе,
эту зависимость можно выразить булевой функцией фиксированного размера,
которая записывается в виде схемы. Таким образом, можно вычислить элементы всех ячеек протокола, а затем проверить, что в последней строке есть
принимающее состояние.

\item $\p \neq \ppoly$

Например, в $\ppoly$ будет лежать любой унарный язык $U$,
т.е. язык, содержащий только слова вида $1^k$. Действительно, для каждой длины $U \cap \{ 0, 1\}^k$ будет либо пустым множеством, либо состоять только из $1^k$. В первом случае язык распознаётся тождественно ложной схемой, во втором —
конъюнкцией.

А унарные языки могут быть даже неразрешимыми, например \newline \textbf{UHALT} = $\{1^k \ | \text{ машина Тьюринга }  M_k  \text{ останавливается на входе }k\}$. 

Также они могут быть разрешимыми, но не лежать в $\p$. Например, можно взять любой язык $A \in \text{\textbf{EEXP}} \backslash \text{\textbf{EXP}}$ (такой есть по теореме об иерархии) и рассмотреть $\{1^k | k \in A \}$. Если его можно распознать за полиномиальное время, то $A$ можно распознать за экспоненциальное время, что противоречит выбору $A$.
\end{itemize}
\end{proof}

\subsubsection{$\mathbf{NC}^d \subset \mathbf{AC}^d \subset \mathbf{NC}^{d+1}$, $\mathbf{NC}^0 \subset \mathbf{AC}^0$}

\begin{statement}
    Для всех $d$ выполнено $\mathbf{NC}^d \subset \mathbf{AC}^d \subset \mathbf{NC}^{d+1}$.
\end{statement}

\begin{proof}
    8.19 у Мусатова, надо вставить % надо бы...
\end{proof}

\begin{statement}
    $\mathbf{NC}^0 \neq \mathbf{AC}^0$.
\end{statement}

\begin{proof}
    В классе $\mathbf{NC}^0$
лежат языки, которые вычисляются схемами
некоторой константной глубины $c$. Поскольку у каждого элемента не больше
двух входов, число изначальных входов, от которых зависит значение, будет
не больше $2^c$. Поэтому такой схемой нельзя вычислить функцию, которая зависит от всех входов, например, их конъюнкцию. В то же время конъюнкция
тривиальным образом лежит в $\mathbf{AC}^0$, откуда $\mathbf{NC}^0 \neq \mathbf{AC}^0$.
\end{proof}

\subsubsection{Вероятностная машина Тьюринга.}

Стоит задача распознавания. Применяем вероятностный алгоритм, который может
с некоторой вероятностью ошибиться. Цель разработки алгоритмов состоит в
том, чтобы вероятность ошибки была малой, но при этом ответ получался достаточно быстро и с помощью небольшого числа случайных битов. 

\begin{statement}
    Под вероятностной машиной Тьюринга мы будем понимать детерминированную машину Тьюринга $M$ с двумя аргументами $x$ (собственно аргумент) и $r$ (случайные биты), где длина $r$ есть некоторая функция от длины $x$. Результатом работы $M$ на входе $x$ будет вероятностное распределение, индуцированное данным $x$ и равномерным распределением на всех значениях $r$. Временем работы $M$ на данном $x$ будем считать максимальное время работы $M(x, r)$ для всех $r$ указанной длины. Так же определяется и использованная память.
\end{statement}

Удобно иметь в виду следующие эквивалентные вариации определения:

\begin{itemize}
\item У машины есть специальная (бесконечная) лента только для чтения, полностью заполненная независимыми равномерно распределёнными случайными битами. Время работы считается как функция от длины основного
входа $x$. Это определение эквивалентно исходному, потому что машина
всё равно прочтёт только конечную область случайных битов.
\item У машины есть две функции перехода, каждый раз машина выбирает
одну из них равновероятно и независимо от предыдущих выборов. Это
определение эквивалентно исходному, поскольку машина может сначала
сгенерировать все случайные биты, а потом работать как раньше (здесь
две функции перехода начнут совпадать).
\item Наконец, у машины может быть специальное состояние, после прихода
в которое машина получает новый случайный бит (например, равновероятно и независимо от предыстории переходя в одно из ещё двух специальных состояний). Это определение потенциально не эквивалентно с
точки зрения использованной памяти: для хранения старых случайных
битов нужна дополнительная память, а на отдельной ленте они хранятся
автоматически.
\end{itemize}

Пример модели в билете 16.

\subsubsection{Классы $\mathbf{BPP}$, $\mathbf{RP}$, $\mathbf{coRP}$, $\mathbf{ZPP}$, $\mathbf{PP}$}.

\begin{definition}
    Классом \textbf{BPP} называется класс языков $A$, для которых существует полиномиальный в худшем случае вероятностный алгоритм $V$, такой что:

    \begin{itemize}
    \item  если $x \in A$, то Pr$_r [V(x, r) = 1] \geq \frac{2}{3}$
    \item  если $x \not \in A$, то Pr$_r [V(x, r) = 1] \leq \frac{1}{3}$.
    \end{itemize}

Иными словами, вероятности ошибок и первого, и второго рода должны
быть не больше $\frac{1}{3}$.

Если поставить $> \frac{1}{2}$ и $< \frac{1}{2}$ вместо $\geq \frac{2}{3}$ и $\leq \frac{1}{3}$
соответственно, получится определение
класса \textbf{PP} (probabilistic polynomial).
\end{definition}

\begin{definition}
     Классом \textbf{RP} называется класс языков $A$, для которых существует полиномиальный в худшем случае вероятностный алгоритм $V$, такой что:

    \begin{itemize}
    \item  если $x \in A$, то Pr$_r [V(x, r) = 1] \geq \frac{1}{2}$
    \item  если $x \not \in A$, то Pr$_r [V(x, r) = 1] = 0$.
    \end{itemize}
\end{definition}

\begin{definition}
    Классом \textbf{coRP} называется класс языков $A$, для которых существует полиномиальный в худшем случае вероятностный алгоритм $V$ , такой что:

    \begin{itemize}
    \item  если $x \in A$, то Pr$_r [V(x, r) = 1] = 1$
    \item  если $x \not \in A$, то Pr$_r [V(x, r) = 1] \leq \frac{1}{2}$.
    \end{itemize}

    Таким образом, для класса \textbf{RP} (randomized polynomial) равна нулю вероятность ошибки первого рода, а для \textbf{coRP} — второго.
    
\end{definition}

\begin{definition}
    Классом \textbf{ZPP} называется множество языков $A$, для которых существует вероятностный алгоритм $V$, такой что $x \in A$ при всех $r$ эквивалентно $V(x, r) = 1$, а для каждого $x$ ожидаемое по $r$ время работы полиномиально.
\end{definition}

Ясно, что $\p \subset \mathbf{ZPP}$, $\p \subset \mathbf{RP} \subset \mathbf{BPP}$ и $\p \subset \mathbf{coRP} \subset \mathbf{BPP}$. Также ясно,
что если $A \in \mathbf{BPP}$, то и $\overline{A} \in \mathbf{BPP}$, ведь определение симметрично. Поэтому $\mathbf{BPP} = \mathbf{coBPP}$. Вот ещё несколько несложных соотношений:

\begin{statement}
    $\mathbf{RP} \subset \mathbf{NP}$.
\end{statement}

\begin{proof}
 Действительно, любое значение $r$, при котором $V(x, r) = 1$, будет доказательством того, что $x \in A$. Можно сказать, что \np тоже определяется вероятностным способом: если $x \in A$, то Pr$_r [V(x, r) = 1] > 0$, а если $x \not \in A$,
то Pr$_r[V(x, r) = 1] = 0$.    
\end{proof}

\begin{statement}
    $\mathbf{PP} \subset \psp.$
\end{statement}

\begin{proof}
Действительно, долю $r$ можно найти точно и сравнить её с $\frac{1}{2}$.
\end{proof}

\begin{statement}
    $\mathbf{BPP} \subset \mathbf{PP}$ - очевидно по определению (в $\mathbf{PP}$ более сильные ограничения)
\end{statement}

